\section{Lecture 14: Neural Networks as Compositional Models}

\subsection{Motivation}

In the previous lecture, we argued that many real-world functions exhibit compositional structure and that learning representations is essential.

\medskip

\noindent
\textbf{Key question.}  
How can we construct models that explicitly capture compositional structure?

\medskip

\noindent
\textbf{Answer.}  
Neural networks implement function composition through layered architectures.

\subsection{Layered Architectures}

A neural network is defined as a composition of layers:
\[
f(x) = f_L \circ f_{L-1} \circ \cdots \circ f_1(x).
\]

\medskip

\noindent
Each layer has the form:
\[
f_\ell(h) = \sigma(W_\ell h + b_\ell),
\]
where:
\begin{itemize}
    \item $W_\ell$ is a weight matrix,
    \item $b_\ell$ is a bias vector,
    \item $\sigma$ is a nonlinear activation function.
\end{itemize}

\medskip

\noindent
\textbf{Forward pass.}

\[
h_1 = \sigma(W_1 x + b_1),
\]
\[
h_2 = \sigma(W_2 h_1 + b_2),
\]
\[
\vdots
\]
\[
f(x) = W_L h_{L-1} + b_L.
\]

\medskip

\noindent
\textbf{Interpretation.}
\begin{itemize}
    \item Each layer transforms its input into a new representation
    \item Representations become progressively more abstract
\end{itemize}

\medskip

\noindent
\textbf{Insight.}  
Neural networks are structured compositions of simple functions.

\subsection{Expressivity and Depth}

\textbf{Key question.}  
Why use multiple layers instead of a single large layer?

\medskip

\noindent
\textbf{Observation.}
\begin{itemize}
    \item Shallow networks can approximate many functions
    \item But may require exponentially many units
\end{itemize}

\medskip

\noindent
\textbf{Depth advantage.}
\begin{itemize}
    \item Deep networks can represent certain functions more efficiently
    \item Composition allows reuse of intermediate computations
\end{itemize}

\medskip

\noindent
\textbf{Example intuition.}
\begin{itemize}
    \item Repeated composition builds hierarchical features
\end{itemize}

\medskip

\noindent
\textbf{Tradeoff.}
\begin{itemize}
    \item Depth increases expressivity
    \item Also increases optimization complexity
\end{itemize}

\medskip

\noindent
\textbf{Insight.}  
Depth enables efficient representation of complex functions.

\subsection{Universal Approximation (Insight)}

\textbf{Statement (informal).}  
A neural network with a single hidden layer and sufficient width can approximate any continuous function on a compact domain.

\medskip

\noindent
\textbf{Implications.}
\begin{itemize}
    \item Neural networks are highly expressive
    \item Depth is not required for expressivity in principle
\end{itemize}

\medskip

\noindent
\textbf{Limitation.}
\begin{itemize}
    \item Approximation may require very large width
\end{itemize}

\medskip

\noindent
\textbf{Key insight.}
\begin{itemize}
    \item Depth is about \emph{efficiency}, not just expressivity
\end{itemize}

\subsection{Representation Learning}

Neural networks learn feature maps automatically:
\[
\phi_\theta(x) = h_{L-1}.
\]

\medskip

\noindent
\textbf{Model.}
\[
f(x) = w^\top \phi_\theta(x).
\]

\medskip

\noindent
\textbf{Interpretation.}
\begin{itemize}
    \item Early layers learn low-level features
    \item Later layers learn high-level abstractions
\end{itemize}

\medskip

\noindent
\textbf{Contrast with classical methods.}
\begin{itemize}
    \item Classical ML:
    \begin{itemize}
        \item Fixed features
    \end{itemize}
    \item Neural networks:
    \begin{itemize}
        \item Learned features
    \end{itemize}
\end{itemize}

\medskip

\noindent
\textbf{Insight.}  
Representation learning is a central advantage of deep learning.

\subsection{A Unifying Perspective}

Neural networks unify key ideas from earlier lectures:

\begin{itemize}
    \item \textbf{Function approximation:} large hypothesis spaces
    \item \textbf{Optimization:} gradient-based training
    \item \textbf{Representation:} learned feature maps
    \item \textbf{Composition:} hierarchical structure
\end{itemize}

\medskip

\noindent
\textbf{Core idea.}  
A neural network is a parameterized composition of functions trained end-to-end.

\medskip

\noindent
\textbf{Key insight.}  
The power of neural networks lies in combining representation learning with compositional structure.

\subsection{Outlook}

In this lecture, we introduced:

\begin{itemize}
    \item layered neural network architectures,
    \item expressivity and the role of depth,
    \item the universal approximation principle,
    \item representation learning.
\end{itemize}

\medskip

\noindent
To train neural networks, we need efficient methods to compute gradients.

\medskip

\noindent
In the next lecture, we study backpropagation and automatic differentiation, which enable scalable training of deep models.

\medskip

\noindent
\textbf{Guiding principle.}  
Deep models learn powerful representations through composition and optimization.